
\documentclass[12pt]{article}
%%%%%%%%%%%%%%%%%%%%%%%%%%%%%%%%%%%%%%%%%%%%%%%%%%%%%%%%%%%%%%%%%%%%%%%%%%%%%%%%%%%%%%%%%%%%%%%%%%%%%%%%%%%%%%%%%%%%%%%%%%%%%%%%%%%%%%%%%%%%%%%%%%%%%%%%%%%%%%%%%%%%%%%%%%%%%%%%%%%%%%%%%%%%%%%%%%%%%%%%%%%%%%%%%%%%%%%%%%%%%%%%%%%%%%%%%%%%%%%%%%%%%%%%%%%%
\usepackage{amsmath}
\usepackage[compat2]{geometry}

\setcounter{MaxMatrixCols}{10}
%TCIDATA{TCIstyle=LaTeX article (bright).cst}

%TCIDATA{OutputFilter=LATEX.DLL}
%TCIDATA{Version=5.50.0.2953}
%TCIDATA{<META NAME="SaveForMode" CONTENT="1">}
%TCIDATA{BibliographyScheme=Manual}
%TCIDATA{Created=Monday, May 23, 2011 14:12:05}
%TCIDATA{LastRevised=Monday, November 27, 2023 12:13:20}
%TCIDATA{<META NAME="GraphicsSave" CONTENT="32">}
%TCIDATA{<META NAME="DocumentShell" CONTENT="Exams and Syllabi\SW\Assignment">}

\setlength{\topmargin}{-1.0in}
\setlength{\textheight}{9.25in}
\setlength{\oddsidemargin}{0.0in}
\setlength{\evensidemargin}{0.0in}
\setlength{\textwidth}{6.5in}
\def\labelenumi{\arabic{enumi}.}
\def\theenumi{\arabic{enumi}}
\def\labelenumii{(\alph{enumii})}
\def\theenumii{\alph{enumii}}
\def\p@enumii{\theenumi.}
\def\labelenumiii{\arabic{enumiii}.}
\def\theenumiii{\arabic{enumiii}}
\def\p@enumiii{(\theenumi)(\theenumii)}
\def\labelenumiv{\arabic{enumiv}.}
\def\theenumiv{\arabic{enumiv}}
\def\p@enumiv{\p@enumiii.\theenumiii}
\pagestyle{plain}
\setcounter{secnumdepth}{0}
\parindent=0pt
\input{tcilatex}
\begin{document}


\begin{center}
\begin{equation*}
\FRAME{itbpF}{1.3214in}{0.9496in}{0in}{}{}{Figure}{\special{language
"Scientific Word";type "GRAPHIC";maintain-aspect-ratio TRUE;display
"USEDEF";valid_file "T";width 1.3214in;height 0.9496in;depth
0in;original-width 1.3552in;original-height 0.9669in;cropleft "0";croptop
"1";cropright "1";cropbottom "0";tempfilename
'R2D0SD00.bmp';tempfile-properties "XPR";}}
\end{equation*}%
\textbf{Electromagnetismo}

\emph{Gu\'{\i}a N}$%
%TCIMACRO{\U{b0}}%
%BeginExpansion
{{}^\circ}%
%EndExpansion
8$\emph{\ - Fuerza Magn\'{e}tica sobre cargas y corrientes}

\textit{Profesor: Arturo\ Fern\'{a}ndez P\'{e}rez}
\end{center}

\begin{enumerate}
\item Un prot\'{o}n que se mueve a $4\times 10^{6}$ [m/s] a trav\'{e}s de
campo magn\'{e}tico de $1.7$ [T] experimenta una fuerza magn\'{e}tica de
magnitud $8.2\times 10^{-13}$ [N] \textquestiondown Cu\'{a}l es el \'{a}%
ngulo entre la velocidad y el campo? \textbf{R:} $\theta =48,91%
%TCIMACRO{\U{b0}}%
%BeginExpansion
{{}^\circ}%
%EndExpansion
$

\item Un prot\'{o}n se mueve perpendicular a un campo magn\'{e}tico uniforme
con una rapidez de $1\times 10^{7}$ [m/s] y experimenta una aceleraci\'{o}n
de $2\times 10^{13}$ [m/s$^{2}$] en la direcci\'{o}n $+x$ cuando su
velocidad est\'{a} en la direcci\'{o}n $+z$. Determine la magnitud y direcci%
\'{o}n del campo magn\'{e}tico. \textbf{R: }$\mathbf{B}=-0.02\mathbf{j}$ [T]

\item Una part\'{\i}cula alfa de carga $q=3.2\times 10^{-19}$ [C] se mueve
con una velocidad $\mathbf{v}=70\mathbf{i}-40\mathbf{j}+100\mathbf{k}$ [m/s]
dentro de una regi\'{o}n que posee un campo magn\'{e}tico uniforme $\mathbf{B%
}=3\mathbf{i}+20\mathbf{j}-5\mathbf{k}$ [T] y un campo el\'{e}ctrico
uniforme $\mathbf{E}=40\mathbf{i}+95\mathbf{j}+23\mathbf{k}$ [V/m].
Determine la fuerza de Lorentz que se induce sobre la part\'{\i}cula alfa. 
\textbf{R: }$\mathbf{F}=-5,632\times 10^{-16}\mathbf{i}+2,384\times 10^{-16}%
\mathbf{j}+4,938\times 10^{-16}\mathbf{k}$ [N]

\item Un electr\'{o}n experimenta una fuerza magn\'{e}tica cuya magnitud es
de $4.6\times 10^{-15}$ [N] cuando se desplaza a un \'{a}ngulo de 60%
%TCIMACRO{\U{ba} }%
%BeginExpansion
${{}^o}$
%EndExpansion
con respecto a un campo magn\'{e}tico con un magnitud de $3.5\times 10^{-3}$
[T].\ Determine la rapidez del electr\'{o}n. \textbf{R: }$v=9,48\times
10^{6} $ [m/s]

\item Un electr\'{o}n se acelera desde el reposo por una diferencia de
potencial de 350 [V]. Luego, entra a una regi\'{o}n de campo magn\'{e}tico
uniforme de 200 [mT] de magnitud, con el que su velocidad forma un \'{a}%
ngulo recto. Determine la rapidez del electr\'{o}n y el radio de la
trayectoria en el campo magn\'{e}tico. \textbf{R: }$v=1,10\times 10^{7}$
[m/s] , $R=3,12\times 10^{-4}$ [m]

\item Un haz de electrones entran en una regi\'{o}n con campos el\'{e}ctrico
y magn\'{e}tico uniformes de magnitudes $1.56\times 10^{4}$ [V/m] y $%
4.62\times 10^{-3}$ [T], respectivamente. Si la velocidad y los campos son
perpendiculares entre s\'{\i} y los electrones no son deflectados, determine
su rapidez y dibuje en un esquema las direcciones de cada uno de los
vectores $\mathbf{E}$, $\mathbf{B}$\textbf{\ }y $\mathbf{v}$. \textbf{R: }$%
v=3,38\times 10^{6}$ [m/s]

\item El cubo de la Fig. 1, de $55$ [cm] de lado, est\'{a} en un campo magn%
\'{e}tico uniforme de magnitud $4$ [T] paralelo al eje $y$ positivo. Cuatro
segmentos rectos $ab$, $bc,$ $cd$ y $da$ forman un lazo cerrado que conduce
una corriente $I=4.3$ [A], en la direcci\'{o}n que muestra la figura.
Determina la fuerza magn\'{e}tica que act\'{u}a sobre cada uno de los
segmentos y la fuerza neta que act\'{u}a sobre el alambre. \textbf{R: }$%
\mathbf{F}_{ab}=0$ [N]$,$ $\mathbf{F}_{bc}=-9,46\mathbf{i}$ [N]$,$ $\mathbf{F%
}_{cd}=-9,46\mathbf{k}$ [N], $\mathbf{F}_{da}=9,46\left( \mathbf{i+k}\right) 
$ [N], $\mathbf{F}=0$ [N] 
\begin{equation*}
\FRAME{itbpFU}{2.1482in}{1.7469in}{0in}{\Qcb{Figura 1}}{}{Figure}{\special%
{language "Scientific Word";type "GRAPHIC";maintain-aspect-ratio
TRUE;display "USEDEF";valid_file "T";width 2.1482in;height 1.7469in;depth
0in;original-width 2.6463in;original-height 2.1465in;cropleft "0";croptop
"1";cropright "1";cropbottom "0";tempfilename
'R2D0SD01.bmp';tempfile-properties "XPR";}}
\end{equation*}

\item Un cable transporta una corriente de 50~[A] de Oeste a Este, en una
regi\'{o}n donde hay un campo magn\'{e}tico uniforme en direcci\'{o}n 45%
%TCIMACRO{\U{ba} }%
%BeginExpansion
${{}^o}$
%EndExpansion
al Norte del Este, que tiene una magnitud de 2 [T], como se muestra en la
Fig. 2. Determine la fuerza magn\'{e}tica ejercida sobre una secci\'{o}n de
1 [m] de cable. Suponiendo que la gravedad apunta en direcci\'{o}n opuesta a
la fuerza resultante \textquestiondown Qu\'{e} masa debe tener el cable de
cobre para que experimente levitaci\'{o}n magn\'{e}tica? (Peso equilibrado
con la fuerza magn\'{e}tica) \textbf{R:} $\mathbf{F}=70,71\mathbf{k}$ [N], $%
m=7,21$ [kg]%
\begin{equation*}
\FRAME{itbpFU}{2.6714in}{1.8282in}{0in}{\Qcb{Figura 2}}{}{Figure}{\special%
{language "Scientific Word";type "GRAPHIC";maintain-aspect-ratio
TRUE;display "USEDEF";valid_file "T";width 2.6714in;height 1.8282in;depth
0in;original-width 3.7161in;original-height 2.5356in;cropleft "0";croptop
"1";cropright "1";cropbottom "0";tempfilename
'R2D0SD02.bmp';tempfile-properties "XPR";}}
\end{equation*}

\item Un alambre conduce una corriente de 1.2 [A] hacia el centro de la
Tierra (direcci\'{o}n perpendicular al eje Norte-Sur) en una regi\'{o}n de
un campo magn\'{e}tico uniforme de magnitud 0.588 [T] en el plano de la
Tierra. \textquestiondown Cu\'{a}l es la direcci\'{o}n y magnitud de la
fuerza por unidad de longitud, si el campo magn\'{e}tico apunta

\begin{description}
\item[a.] Hacia el Este? \textbf{R: }$0,7056$ [N] en direcci\'{o}n Sur.

\item[b.] Hacia el Sur? \ \textbf{R: }$0,7056$ [N] en direcci\'{o}n Oeste.

\item[c.] 30%
%TCIMACRO{\U{ba} }%
%BeginExpansion
${{}^o}$
%EndExpansion
al Sur del Este? \textbf{R: }$0,7056$ [N] a 30$%
%TCIMACRO{\U{b0}}%
%BeginExpansion
{{}^\circ}%
%EndExpansion
$ al Oeste del sur.
\end{description}

\item Una espira rectangular de lados 22 y 35 [cm] conduce una corriente $%
I=1.4$ [A]. Si la espira est\'{a} dentro de una regi\'{o}n de campo magn\'{e}%
tico uniforme $B=1.5$ [T], determine la fuerza neta sobre la espira y el
torque que se ejerce alrededor del eje ("axis") de rotaci\'{o}n. (Fig. 4) 
\textbf{R: }$\mathbf{F=0}$ [N]\textbf{, }$\mathbf{\tau =0}$ [Nm]%
\begin{equation*}
\FRAME{itbpFU}{2.0634in}{1.2211in}{0in}{\Qcb{Figura 4}}{}{Figure}{\special%
{language "Scientific Word";type "GRAPHIC";maintain-aspect-ratio
TRUE;display "USEDEF";valid_file "T";width 2.0634in;height 1.2211in;depth
0in;original-width 3.2093in;original-height 1.8896in;cropleft "0";croptop
"1";cropright "1";cropbottom "0";tempfilename
'R2D0SD03.wmf';tempfile-properties "XPR";}}
\end{equation*}

\item El lazo rectangular pivotea en el eje $y$, como muestra la Figura 5.
Por \'{e}l circula una corriente $I=15$ [A]. Si este lazo est\'{a} en una
regi\'{o}n con un campo magn\'{e}tico uniforme de 0.48 [T] en la direcci\'{o}%
n $x+$ , determine el torque necesario para mantener fijo al lazo en la
posici\'{o}n que se muestra en la Fig. 5. Luego, determine la magnitud de la
fuerza que debe aplicarse perpendicular al lado opuesto al eje, de manera de
mantenerlo en esa posici\'{o}n. \textbf{R: }$\mathbf{\tau }=0,03\mathbf{j}$
[Nm] , $F_{ext}=0,5$ [N].%
\begin{equation*}
\FRAME{itbpFU}{1.4468in}{1.7244in}{0in}{\Qcb{Figura 5}}{}{Figure}{\special%
{language "Scientific Word";type "GRAPHIC";maintain-aspect-ratio
TRUE;display "USEDEF";valid_file "T";width 1.4468in;height 1.7244in;depth
0in;original-width 3.3736in;original-height 4.0274in;cropleft "0";croptop
"1";cropright "1";cropbottom "0";tempfilename
'R2D0SD04.wmf';tempfile-properties "XPR";}}
\end{equation*}

\item Un cable de masa $m=13$ [g] y largo $L=62$ [cm] se encuentra
suspendido por 2 cuerdas flexibles, como se muestra en la Fig. 6, en un
campo magn\'{e}tico uniforme de magnitud 0.44 [T] \textquestiondown Cu\'{a}l
es la magnitud y direcci\'{o}n de la corriente requerida para que la tensi%
\'{o}n sobre las cuerdas sea cero? \textbf{R: }$I=0,47$ [A] en sentido
anti-horario 
\begin{equation*}
\FRAME{itbpFU}{1.516in}{0.9124in}{0in}{\Qcb{Figura 6}}{}{Figure}{\special%
{language "Scientific Word";type "GRAPHIC";maintain-aspect-ratio
TRUE;display "USEDEF";valid_file "T";width 1.516in;height 0.9124in;depth
0in;original-width 3.0398in;original-height 1.8196in;cropleft "0";croptop
"1";cropright "1";cropbottom "0";tempfilename
'R2D0SD05.wmf';tempfile-properties "XPR";}}
\end{equation*}

\item Un solenoide de $50$ vueltas y radio $R=5,0$ [cm], se sit\'{u}a en una
regi\'{o}n de campo magn\'{e}tico uniforme de magnitud $B=0,6$ [T], como
muestra la Fig. 7. Si por el solenoide circula una corriente $I=2,0$ [A],
determine la magnitud del momento magn\'{e}tico del solenoide $\mu $, la
magnitud del torque ejercido sobre \'{e}l $\tau $ y el sentido de rotaci\'{o}%
n del mismo. \textbf{R: }$\mu =0,79$ [Am$^{2}$], $\tau =0,30$ [Nm$^{2}$],
rota en sentido horario. 
\begin{equation*}
\FRAME{itbpFU}{2.0686in}{1.7608in}{0in}{\Qcb{Figura 7}}{}{Figure}{\special%
{language "Scientific Word";type "GRAPHIC";maintain-aspect-ratio
TRUE;display "USEDEF";valid_file "T";width 2.0686in;height 1.7608in;depth
0in;original-width 2.3748in;original-height 2.0176in;cropleft "0";croptop
"1";cropright "1";cropbottom "0";tempfilename
'R2PVIR00.wmf';tempfile-properties "XPR";}}
\end{equation*}
\end{enumerate}

\end{document}
